% !TeX TS-program = xelatex
% 虚构演示简历 — 仅用于 hellojobs 测试数据
\documentclass{resume-photo}
\usepackage{xeCJK}
\setCJKmainfont{Noto Serif CJK SC}
\usepackage{fontspec}
\usepackage{enumitem}
\setlist[itemize]{itemsep=0pt, topsep=1pt, parsep=0pt, partopsep=0pt}

\ResumeName{孙悦}

\begin{document}

\ResumeContacts{
  138-0007-0007,%
  \ResumeUrl{mailto:sunyue@buaa.edu.cn}{sunyue@buaa.edu.cn},%
  \textnormal{北京航空航天大学 | 软件工程 · 硕士 | 2023—2026}%
}

\ResumeTitle

\noindent iOS 与 SwiftUI，熟悉 Combine、Core Data 与网络层封装；实习期间负责金融 App 安全键盘与证书钉扎改造。注重可观测性与文档沉淀，习惯在评审阶段拆解风险；参与线上复盘与跨团队联调，英语 CET-6，可阅读官方文档与 RFC（演示数据）。

\section{教育背景}
\ResumeItem
{北京航空航天大学}
[\textnormal{软件学院 | 软件工程 | 硕士}]
[2023.09—2026.06(预计)]

\begin{itemize}
  \item 移动与嵌入式软件方向。
\end{itemize}


\ResumeItem
{国际交流（演示）}
[\textnormal{暑期学校 / 线上联合项目}]
[2023—2024]

\begin{itemize}
  \item 参加海外高校暑期学校（在线），完成 Distributed Systems 课程 Project 3 个。
  \item 与海外小组合作完成跨时区项目，使用 Slack/Discord 进行异步协作与代码评审。
  \item 撰写全英文技术报告 8 页，获得课程组 Best Presentation 提名（测试数据）。
  \item 口语与写作训练：技术博客英文化 5 篇，累计阅读量 1.2 万（演示）。
\end{itemize}



\section{专业技能}
\begin{itemize}
  \item 掌握 Swift、UIKit/SwiftUI、Xcode Instruments。
  \item 熟悉 HTTPS 证书钉扎、Keychain 与防截屏策略。
  \item 了解 Objective-C 互操作与 SPM 模块化。
  \item 熟悉 Linux 日常运维与 Shell，具备日志检索与 CPU/内存突增初步定位能力。
  \item 掌握 Git / CI 与 Code Review，了解 Docker 与 Kubernetes 基本编排与滚动发布。
  \item 了解 OAuth2 / JWT、接口幂等与 REST 规范，具备单元测试与集成测试习惯（JUnit/pytest 等）。
  \item 能阅读英文技术文档，具备跨团队沟通与需求澄清能力（测试简历）。
\end{itemize}

\section{实习经历}
\ResumeItem
{东方财富}
[\textnormal{iOS 开发实习生}]
[2024.07—2024.12]

\begin{itemize}
  \item \textbf{职责}: 交易模块崩溃率治理与启动框架拆分。
  \item \textbf{成果}: 崩溃率从 0.35\% 降至 0.12\%（周活维度）。
\end{itemize}


\ResumeItem
{海康威视 | 行业应用}
[\textnormal{嵌入式 / 后端实习生}]
[2023.09—2024.02]

\begin{itemize}
  \item \textbf{设备接入}: 参与国标设备接入网关协议适配与异常码映射。
  \item \textbf{可靠性}: 实现断线重连与心跳退避策略，掉线恢复时间缩短 40\%。
  \item \textbf{日志}: 统一日志格式与 traceId 透传，问题定位时间减半。
  \item \textbf{客户}: 支持 2 次驻场联调，收集需求并转化为可验收用例。
  \item \textbf{规范}: 熟悉编码规范与静态检查门禁，CI 全绿方可合并。
\end{itemize}



\section{项目经历}
\ResumeItem
{SwiftUI 股票分时图组件}
[\textnormal{课程设计}]
[2024.03—2024.06]

\begin{itemize}
  \item Canvas 绘制 + 手势缩放，帧率稳定在 55+ FPS（iPhone 12）。
\end{itemize}


\ResumeItem
{灰度发布与流量染色}
[\textnormal{网关 + 标签路由（演示）}]
[2024.02—2024.06]

\begin{itemize}
  \item \textbf{能力}: 基于用户 ID / 地域 / 白名单实现多维度灰度与染色。
  \item \textbf{实现}: 扩展网关插件，透传染色 Header 至下游 RPC Metadata。
  \item \textbf{验证}: 与 QA 共建自动化回归集，覆盖 200+ 核心用例。
  \item \textbf{效果}: 月均生产配置事故从 3 起降至 0.5 起（演示数据）。
\end{itemize}



\section{个人荣誉}
\begin{itemize}
  \item 中国大学生服务外包创新创业大赛 国家三等奖
  \item 北航 学习优秀奖学金
  \item 全国计算机等级考试四级（网络工程师）（演示）
  \item 校级「优秀学生干部 / 优秀团员」或技术社团骨干（综合测评前 15\%）
  \item 开源贡献：向知名项目提交被合并的 PR（文档/测试/feature）（演示）
\end{itemize}

\end{document}
